% Part: incompleteness
% Chapter: incompleteness-provability
% Section: introduction

\documentclass[../../../include/open-logic-section]{subfiles}

\begin{document}

\olfileid{inc}{inp}{int}

\olsection{تعارف}

ہلبرٹ دا خیال سی کہ قدرتی اعداد ورگی کسے ریاضیاتی ساخت دے
مسلّمیاں دا نظام اوہدوں ای کافی اے جدوں اوہدے توں اوس ساخت
بارے ہر سچا بیان اخذ ہو سکے۔ پچھوں رسمی نظاماں دے ذریعے اخذ وچ
اوہدی دلچسپی توں لگدا اے کہ اوہ ایہ ضمانت چاہندا سی کہ مثلاً
قدرتی اعداد بارے دلیل دین لئی ورتے رسمی نظام بے تضاد وی ہون
تے \emph{مکمل} وی۔ یعنی اوہناں دی زبان دا ہر بیان یا اوہدی نفی
!!{derivable} ہووے۔ گوڈل دا پہلا نامکملیت قضیہ وکھاؤندا اے کہ
مسلّمیاں دا ایہو جیہا کوئی نظام موجود نہیں۔ حساب دا کوئی رسمی
نظام مکمل، بے تضاد تے !!{axiomatizable} تینوں اکّو ویلے نہیں ہو سکدا۔
بلکہ کوئی وی “کافی طاقتور” بے تضاد تے !!{axiomatizable}
ریاضیاتی نظریہ مکمل نہیں۔

ہلبرٹ دا اک ہور، ودھ اہم مقصد، جدید (“کلاسیکی”) ریاضی دا
جواز قائم کرن دے پروگرام دا مرکز سی۔ اوہ کلاسیکی استدلال دی
نمائندگی کرن والے رسمی نظاماں دی بے تضادی دے ثبوت محدود
طریقیاں نال لبھنا چاہندا سی۔ ایس لئی گوڈل دا دوجا نامکملیت
قضیہ ہلبرٹ دے پروگرام لئی کہیں وڈا دھچکا سی۔ پہلے قضیے وانگ
ایہنوں وی موٹے طور اُتے بیان کیتا جا سکدا اے۔ حساب دا کوئی
کافی طاقتور نظریہ اپنی بے تضادی ثابت نہیں کر سکدا۔
“کافی طاقتور” آکھدیاں سانوں~$\Th{Q}$ توں کجھ ودھ طاقت
مراد لینی پئے گی۔

گوڈل دے نامکملیت قضیے دے اصل ثبوت دا خیال ایپیمنیدس دے
پیراڈوکس وچ لبھدا اے۔ کریٹ دا ایپیمنیدس آکھدا سی کہ کریٹ دے
سارے لوک جھوٹ بولدے نیں۔ ایس پیراڈوکس دی سدھی صورت ایہ
اے، “ایہ جملہ جھوٹا اے۔” سچائی دی تھاں !!{derivability} رکھ
کے گوڈل نے !!a{sentence} نوں رسمی روپ دتا جیہڑا بالواسطہ
اپنے بارے آکھدا اے کہ اوہ !!{derivable} نہیں۔ جے اوہ
!!{sentence} !!{derivable} ہووے، تاں نظریہ متضاد ہو جائے۔
گوڈل نے ایہ وی وکھایا کہ اوس !!{sentence} دی نفی زیرِ غور
مسلّمیاں دے نظام توں !!{derivable} نہیں۔ (ایس دوجے حصے لئی
گوڈل نوں فرض کرنا پیا کہ نظریہ~$\Th{T}$ “$\omega$-بے تضاد”
اے۔ $\omega$-بے تضادی، عام بے تضادی نال جڑی ہوئی پر ودھ
مضبوط خاصیت اے۔\footnote{یعنی ہر $\omega$-بے تضاد نظریہ
بے تضاد ہوندا اے، پر الٹی شرط لازم نہیں۔} گوڈل توں کجھ سال
پچھوں روسر نے وکھایا کہ نظریے~$\Th{T}$ دی عام بے تضادی
فرض کرنی کافی اے۔)

پہلا سوال ایہ اے کہ اپنے آپ ول حوالہ کرن والا !!a{sentence}
کِویں بنایا جائے۔ ہر !!{formula}~$!A$ لئی جیہڑا~$\Th{Q}$
دی زبان دا اے، منّ لئیے کہ~$\gn{!A}$ اوہ عددی نام اے جیہڑا
~$\Gn{!A}$ دے مطابق اے۔ ذرا غور کرو۔ $!A$،~$\Th{Q}$
دی زبان دا !!a{formula} اے۔ $\Gn{!A}$ اک قدرتی عدد اے،
تے $\gn{!A}$،~$\Th{Q}$ دی زبان دا اک \emph{ٹرم} اے۔
ایس لئی ہر !!{formula}~$!A$، جیہڑا~$\Th{Q}$ دی زبان دا اے،
دا اک \emph{ناں}،~$\gn{!A}$، اے۔ ایہ ناں~$\Th{Q}$ دی زبان
دا ٹرم اے۔ ایس طرح سانوں اوہ ڈھانچا ملدا اے جیہدے وچ
~$\Th{Q}$ دی زبان دے !!{formula}s دوجے !!{formula}s بارے
گل “کر” سکدے نیں۔ اگلے لمے نوں مقرر نقطے دا لمہ آکھدے آں۔

\begin{lem}
فرض کرو $\Th{T}$،~$\Th{Q}$ دی کوئی وی توسیع اے، تے $!B(x)$
کوئی ایہو جیہا !!{formula} اے جیہدے وچ صرف متغیر~$x$ آزاد اے۔
تاں اک !!a{sentence}~$!A$ موجود اے جیہدے لئی
$\Th{T} \Proves!A \liff !B(\gn{!A})$۔
\end{lem}

لمہ آکھدا اے کہ کوئی وی خاصیت $!B(x)$ دتی ہووے تاں اک
!!a{sentence}~$!A$ اے جیہڑا آکھدا اے، “$!B(x)$ میرے بارے
سچ اے۔” $\Th{T}$ ایہ گل “جاندا” اے۔

ایہو جیہا !!a{sentence} اسیں کِویں بناواں؟ ایپیمنیدس دے
پیراڈوکس دی ایہ صورت ویکھو جیہڑی کوائن نے دتی۔
\begin{quote}
“اپنا اقتباس اگے رکھن نال جھوٹ نکلدا اے” اپنا اقتباس اگے
رکھن نال جھوٹ نکلدا اے۔
\end{quote}
ایہ جملہ سِدھا اپنے آپ ول حوالہ نہیں کردا۔ ایہ بس اقتباس
دیاں نشانیاں وچ دتّی عبارت بارے اک نحوی دعویٰ کردا اے۔
ایس لحاظ نال ایہ ہِٹھلے جملیاں ورگا اے۔
\begin{enumerate}
\item “رابرٹ” اک سوہنا ناں اے۔
\item “میں دوڑیا۔” اک چھوٹا جملہ اے۔
\item “تِن لفظ نیں” وچ تِن لفظ نیں۔
\end{enumerate}
پر جے “اپنا اقتباس اگے رکھن نال جھوٹ نکلدا اے” والی عبارت
لئیے تے اوہدے اگے اوسے عبارت دا اقتباس رکھ دیئے، تاں کیہ
ہوندا اے؟ اوہی اصل جملہ بن جاندا اے!{} ایس طرح جملہ اپنے
بارے آکھدا اے کہ ایہ جھوٹا اے۔

\end{document}
